% Options for packages loaded elsewhere
% Options for packages loaded elsewhere
\PassOptionsToPackage{unicode}{hyperref}
\PassOptionsToPackage{hyphens}{url}
\PassOptionsToPackage{dvipsnames,svgnames,x11names}{xcolor}
%
\documentclass[
  letterpaper,
  DIV=11,
  numbers=noendperiod]{scrartcl}
\usepackage{xcolor}
\usepackage[margin=0.7in]{geometry}
\usepackage{amsmath,amssymb}
\setcounter{secnumdepth}{-\maxdimen} % remove section numbering
\usepackage{iftex}
\ifPDFTeX
  \usepackage[T1]{fontenc}
  \usepackage[utf8]{inputenc}
  \usepackage{textcomp} % provide euro and other symbols
\else % if luatex or xetex
  \usepackage{unicode-math} % this also loads fontspec
  \defaultfontfeatures{Scale=MatchLowercase}
  \defaultfontfeatures[\rmfamily]{Ligatures=TeX,Scale=1}
\fi
\usepackage{lmodern}
\ifPDFTeX\else
  % xetex/luatex font selection
\fi
% Use upquote if available, for straight quotes in verbatim environments
\IfFileExists{upquote.sty}{\usepackage{upquote}}{}
\IfFileExists{microtype.sty}{% use microtype if available
  \usepackage[]{microtype}
  \UseMicrotypeSet[protrusion]{basicmath} % disable protrusion for tt fonts
}{}
\makeatletter
\@ifundefined{KOMAClassName}{% if non-KOMA class
  \IfFileExists{parskip.sty}{%
    \usepackage{parskip}
  }{% else
    \setlength{\parindent}{0pt}
    \setlength{\parskip}{6pt plus 2pt minus 1pt}}
}{% if KOMA class
  \KOMAoptions{parskip=half}}
\makeatother
% Make \paragraph and \subparagraph free-standing
\makeatletter
\ifx\paragraph\undefined\else
  \let\oldparagraph\paragraph
  \renewcommand{\paragraph}{
    \@ifstar
      \xxxParagraphStar
      \xxxParagraphNoStar
  }
  \newcommand{\xxxParagraphStar}[1]{\oldparagraph*{#1}\mbox{}}
  \newcommand{\xxxParagraphNoStar}[1]{\oldparagraph{#1}\mbox{}}
\fi
\ifx\subparagraph\undefined\else
  \let\oldsubparagraph\subparagraph
  \renewcommand{\subparagraph}{
    \@ifstar
      \xxxSubParagraphStar
      \xxxSubParagraphNoStar
  }
  \newcommand{\xxxSubParagraphStar}[1]{\oldsubparagraph*{#1}\mbox{}}
  \newcommand{\xxxSubParagraphNoStar}[1]{\oldsubparagraph{#1}\mbox{}}
\fi
\makeatother

\usepackage{color}
\usepackage{fancyvrb}
\newcommand{\VerbBar}{|}
\newcommand{\VERB}{\Verb[commandchars=\\\{\}]}
\DefineVerbatimEnvironment{Highlighting}{Verbatim}{commandchars=\\\{\}}
% Add ',fontsize=\small' for more characters per line
\usepackage{framed}
\definecolor{shadecolor}{RGB}{241,243,245}
\newenvironment{Shaded}{\begin{snugshade}}{\end{snugshade}}
\newcommand{\AlertTok}[1]{\textcolor[rgb]{0.68,0.00,0.00}{#1}}
\newcommand{\AnnotationTok}[1]{\textcolor[rgb]{0.37,0.37,0.37}{#1}}
\newcommand{\AttributeTok}[1]{\textcolor[rgb]{0.40,0.45,0.13}{#1}}
\newcommand{\BaseNTok}[1]{\textcolor[rgb]{0.68,0.00,0.00}{#1}}
\newcommand{\BuiltInTok}[1]{\textcolor[rgb]{0.00,0.23,0.31}{#1}}
\newcommand{\CharTok}[1]{\textcolor[rgb]{0.13,0.47,0.30}{#1}}
\newcommand{\CommentTok}[1]{\textcolor[rgb]{0.37,0.37,0.37}{#1}}
\newcommand{\CommentVarTok}[1]{\textcolor[rgb]{0.37,0.37,0.37}{\textit{#1}}}
\newcommand{\ConstantTok}[1]{\textcolor[rgb]{0.56,0.35,0.01}{#1}}
\newcommand{\ControlFlowTok}[1]{\textcolor[rgb]{0.00,0.23,0.31}{\textbf{#1}}}
\newcommand{\DataTypeTok}[1]{\textcolor[rgb]{0.68,0.00,0.00}{#1}}
\newcommand{\DecValTok}[1]{\textcolor[rgb]{0.68,0.00,0.00}{#1}}
\newcommand{\DocumentationTok}[1]{\textcolor[rgb]{0.37,0.37,0.37}{\textit{#1}}}
\newcommand{\ErrorTok}[1]{\textcolor[rgb]{0.68,0.00,0.00}{#1}}
\newcommand{\ExtensionTok}[1]{\textcolor[rgb]{0.00,0.23,0.31}{#1}}
\newcommand{\FloatTok}[1]{\textcolor[rgb]{0.68,0.00,0.00}{#1}}
\newcommand{\FunctionTok}[1]{\textcolor[rgb]{0.28,0.35,0.67}{#1}}
\newcommand{\ImportTok}[1]{\textcolor[rgb]{0.00,0.46,0.62}{#1}}
\newcommand{\InformationTok}[1]{\textcolor[rgb]{0.37,0.37,0.37}{#1}}
\newcommand{\KeywordTok}[1]{\textcolor[rgb]{0.00,0.23,0.31}{\textbf{#1}}}
\newcommand{\NormalTok}[1]{\textcolor[rgb]{0.00,0.23,0.31}{#1}}
\newcommand{\OperatorTok}[1]{\textcolor[rgb]{0.37,0.37,0.37}{#1}}
\newcommand{\OtherTok}[1]{\textcolor[rgb]{0.00,0.23,0.31}{#1}}
\newcommand{\PreprocessorTok}[1]{\textcolor[rgb]{0.68,0.00,0.00}{#1}}
\newcommand{\RegionMarkerTok}[1]{\textcolor[rgb]{0.00,0.23,0.31}{#1}}
\newcommand{\SpecialCharTok}[1]{\textcolor[rgb]{0.37,0.37,0.37}{#1}}
\newcommand{\SpecialStringTok}[1]{\textcolor[rgb]{0.13,0.47,0.30}{#1}}
\newcommand{\StringTok}[1]{\textcolor[rgb]{0.13,0.47,0.30}{#1}}
\newcommand{\VariableTok}[1]{\textcolor[rgb]{0.07,0.07,0.07}{#1}}
\newcommand{\VerbatimStringTok}[1]{\textcolor[rgb]{0.13,0.47,0.30}{#1}}
\newcommand{\WarningTok}[1]{\textcolor[rgb]{0.37,0.37,0.37}{\textit{#1}}}

\usepackage{longtable,booktabs,array}
\newcounter{none} % for unnumbered tables
\usepackage{calc} % for calculating minipage widths
% Correct order of tables after \paragraph or \subparagraph
\usepackage{etoolbox}
\makeatletter
\patchcmd\longtable{\par}{\if@noskipsec\mbox{}\fi\par}{}{}
\makeatother
% Allow footnotes in longtable head/foot
\IfFileExists{footnotehyper.sty}{\usepackage{footnotehyper}}{\usepackage{footnote}}
\makesavenoteenv{longtable}
\usepackage{graphicx}
\makeatletter
\newsavebox\pandoc@box
\newcommand*\pandocbounded[1]{% scales image to fit in text height/width
  \sbox\pandoc@box{#1}%
  \Gscale@div\@tempa{\textheight}{\dimexpr\ht\pandoc@box+\dp\pandoc@box\relax}%
  \Gscale@div\@tempb{\linewidth}{\wd\pandoc@box}%
  \ifdim\@tempb\p@<\@tempa\p@\let\@tempa\@tempb\fi% select the smaller of both
  \ifdim\@tempa\p@<\p@\scalebox{\@tempa}{\usebox\pandoc@box}%
  \else\usebox{\pandoc@box}%
  \fi%
}
% Set default figure placement to htbp
\def\fps@figure{htbp}
\makeatother





\setlength{\emergencystretch}{3em} % prevent overfull lines

\providecommand{\tightlist}{%
  \setlength{\itemsep}{0pt}\setlength{\parskip}{0pt}}



 


\usepackage{booktabs}
\usepackage{caption}
\usepackage{longtable}
\usepackage{colortbl}
\usepackage{array}
\usepackage{anyfontsize}
\usepackage{multirow}
\usepackage{float}
\floatplacement{figure}{H}
\floatplacement{table}{H}
\usepackage{caption}
\captionsetup[table]{justification=centering}
\usepackage{fvextra}
\DefineVerbatimEnvironment{Highlighting}{Verbatim}{breaklines,commandchars=\\\{\}}
\usepackage{etoolbox}
\BeforeBeginEnvironment{tabular*}{\begin{center}}
\AfterEndEnvironment{tabular*}{\end{center}}
\BeforeBeginEnvironment{Shaded}{\begin{minipage}{\linewidth}}
\AfterEndEnvironment{Shaded}{\end{minipage}\vspace{0.5em}}
\KOMAoption{captions}{tableheading}
\makeatletter
\@ifpackageloaded{caption}{}{\usepackage{caption}}
\AtBeginDocument{%
\ifdefined\contentsname
  \renewcommand*\contentsname{Table of contents}
\else
  \newcommand\contentsname{Table of contents}
\fi
\ifdefined\listfigurename
  \renewcommand*\listfigurename{List of Figures}
\else
  \newcommand\listfigurename{List of Figures}
\fi
\ifdefined\listtablename
  \renewcommand*\listtablename{List of Tables}
\else
  \newcommand\listtablename{List of Tables}
\fi
\ifdefined\figurename
  \renewcommand*\figurename{Figure}
\else
  \newcommand\figurename{Figure}
\fi
\ifdefined\tablename
  \renewcommand*\tablename{Table}
\else
  \newcommand\tablename{Table}
\fi
}
\@ifpackageloaded{float}{}{\usepackage{float}}
\floatstyle{ruled}
\@ifundefined{c@chapter}{\newfloat{codelisting}{h}{lop}}{\newfloat{codelisting}{h}{lop}[chapter]}
\floatname{codelisting}{Listing}
\newcommand*\listoflistings{\listof{codelisting}{List of Listings}}
\makeatother
\makeatletter
\makeatother
\makeatletter
\@ifpackageloaded{caption}{}{\usepackage{caption}}
\@ifpackageloaded{subcaption}{}{\usepackage{subcaption}}
\makeatother
\usepackage{bookmark}
\IfFileExists{xurl.sty}{\usepackage{xurl}}{} % add URL line breaks if available
\urlstyle{same}
% fallback for those not using the hyperref driver hyperxmp:
\makeatletter
\@ifundefined{xmpquote}{\newcommand{\xmpquote}[1]{#1}}{}
\makeatother
\hypersetup{
  pdftitle={ADS-506 Assignment 3.1},
  pdfauthor={Jordan Blackman},
  colorlinks=true,
  linkcolor={blue},
  filecolor={Maroon},
  citecolor={Blue},
  urlcolor={Blue},
  pdfcreator={LaTeX via pandoc}}


\title{ADS-506 Assignment 3.1}
\author{Jordan Blackman}
\date{}
\begin{document}
\maketitle


\begin{Shaded}
\begin{Highlighting}[]
\FunctionTok{library}\NormalTok{(tidyverse)}
\FunctionTok{library}\NormalTok{(fpp3)}
\FunctionTok{library}\NormalTok{(gt)}
\end{Highlighting}
\end{Shaded}

\section{Module 3 Assignment
Exercises}\label{module-3-assignment-exercises}

\subsection{1) Forecasting Australian Wine Sales :
ARIMA}\label{forecasting-australian-wine-sales-arima}

The figure below shows time plots of monthly sales of six types of
Australian wines (red, rose, sweet white, dry white, sparkling, and
fortified) for 1980-1994. Data available in AustralianWines.xls. The
units are thousands of liters.

\pandocbounded{\includegraphics[keepaspectratio]{ADS506-Assignment_3_1_files/figure-pdf/ch5_fig-1.pdf}}

Fortified wine has the largest market share of the six types of wine.
Focus on fortified wine sales alone and produce as accurate a monthly
forecast as possible for the next year using ARIMA models.

\subsubsection{a) Create ACF/PACF plots from a training
set.}\label{a-create-acfpacf-plots-from-a-training-set.}

\begin{itemize}
\tightlist
\item
  Start by partitioning the data using the period until December 1993 as
  the training period.
\item
  Determine the differencing (if any) required to make the data
  stationary (show any test output).
\item
  Plot the ACF/PACF on the stationary time series
\item
  Specify the models suggested by the ACF and PACF plots respectively
\end{itemize}

\begin{Shaded}
\begin{Highlighting}[]
\CommentTok{\# Create a time series object, partition, and check for stationarity}
\NormalTok{aus\_wine\_ts }\OtherTok{\textless{}{-}}\NormalTok{ aus\_wine }\SpecialCharTok{|\textgreater{}} \FunctionTok{select}\NormalTok{(Month, Fortified) }\SpecialCharTok{|\textgreater{}} \FunctionTok{as\_tsibble}\NormalTok{(}\AttributeTok{index =}\NormalTok{ Month)}
\NormalTok{trn }\OtherTok{\textless{}{-}}\NormalTok{ aus\_wine\_ts }\SpecialCharTok{|\textgreater{}} \FunctionTok{filter}\NormalTok{(Month }\SpecialCharTok{\textless{}=} \FunctionTok{yearmonth}\NormalTok{(}\StringTok{"1993 Dec"}\NormalTok{))}
\NormalTok{test }\OtherTok{\textless{}{-}}\NormalTok{ aus\_wine\_ts }\SpecialCharTok{|\textgreater{}} \FunctionTok{filter}\NormalTok{(Month }\SpecialCharTok{\textgreater{}} \FunctionTok{yearmonth}\NormalTok{(}\StringTok{"1993 Dec"}\NormalTok{))}
\NormalTok{ns\_diffs }\OtherTok{\textless{}{-}}\NormalTok{ trn }\SpecialCharTok{|\textgreater{}} \FunctionTok{features}\NormalTok{(Fortified, unitroot\_nsdiffs)}
\NormalTok{trn\_d12 }\OtherTok{\textless{}{-}}\NormalTok{ trn }\SpecialCharTok{|\textgreater{}} \FunctionTok{mutate}\NormalTok{(}\AttributeTok{d12 =} \FunctionTok{difference}\NormalTok{(Fortified, }\DecValTok{12}\NormalTok{))}
\NormalTok{nd\_diffs }\OtherTok{\textless{}{-}}\NormalTok{ trn\_d12 }\SpecialCharTok{|\textgreater{}} \FunctionTok{features}\NormalTok{(d12, unitroot\_ndiffs)}
\NormalTok{kpss }\OtherTok{\textless{}{-}}\NormalTok{ trn\_d12 }\SpecialCharTok{|\textgreater{}} \FunctionTok{features}\NormalTok{(d12, unitroot\_kpss)}

\FunctionTok{tibble}\NormalTok{(}\StringTok{\textasciigrave{}}\AttributeTok{Diagnostic Test}\StringTok{\textasciigrave{}} \OtherTok{=} \FunctionTok{c}\NormalTok{(}\StringTok{"Seasonal Diff (nsdiffs)"}\NormalTok{, }\StringTok{"First Diff (ndiffs)"}\NormalTok{, }\StringTok{"KPSS Test"}\NormalTok{),}
  \StringTok{\textasciigrave{}}\AttributeTok{Series Evaluated}\StringTok{\textasciigrave{}} \OtherTok{=} \FunctionTok{c}\NormalTok{(}\StringTok{"Fortified (Raw)"}\NormalTok{, }\StringTok{"Lag{-}12m diff"}\NormalTok{, }\StringTok{"Lag{-}12m diff"}\NormalTok{),}
  \StringTok{\textasciigrave{}}\AttributeTok{Test Result}\StringTok{\textasciigrave{}} \OtherTok{=} \FunctionTok{c}\NormalTok{(}\FunctionTok{paste0}\NormalTok{(}\StringTok{"D = "}\NormalTok{, ns\_diffs}\SpecialCharTok{$}\NormalTok{nsdiffs), }\FunctionTok{paste0}\NormalTok{(}\StringTok{"d = "}\NormalTok{, nd\_diffs}\SpecialCharTok{$}\NormalTok{ndiffs),}
      \FunctionTok{paste0}\NormalTok{(}\StringTok{"Stat: "}\NormalTok{, }\FunctionTok{round}\NormalTok{(kpss}\SpecialCharTok{$}\NormalTok{kpss\_stat, }\DecValTok{4}\NormalTok{), }\StringTok{" (p = "}\NormalTok{, kpss}\SpecialCharTok{$}\NormalTok{kpss\_pvalue, }\StringTok{")"}\NormalTok{)),}
  \StringTok{\textasciigrave{}}\AttributeTok{Conclusion}\StringTok{\textasciigrave{}} \OtherTok{=} \FunctionTok{c}\NormalTok{(}\StringTok{"1 seasonal diff (D = 1)"}\NormalTok{, }\StringTok{"No regular diff (d = 0)"}\NormalTok{,}
                   \StringTok{"Stationary (Fail to reject H0)"}\NormalTok{)) }\SpecialCharTok{|\textgreater{}}
  \FunctionTok{gt}\NormalTok{() }\SpecialCharTok{|\textgreater{}}
  \FunctionTok{tab\_header}\NormalTok{(}\AttributeTok{title =} \FunctionTok{md}\NormalTok{(}\StringTok{"**Stationarity \& Differencing Diagnostic Tests**"}\NormalTok{),}
             \AttributeTok{subtitle =} \StringTok{"Fortified Wine Sales Training Period (1980–1993)"}\NormalTok{) }\SpecialCharTok{|\textgreater{}}
  \FunctionTok{cols\_align}\NormalTok{(}\AttributeTok{align =} \StringTok{"center"}\NormalTok{, }\AttributeTok{columns =} \FunctionTok{everything}\NormalTok{()) }\SpecialCharTok{|\textgreater{}}
  \FunctionTok{tab\_options}\NormalTok{(}\AttributeTok{data\_row.padding =} \FunctionTok{px}\NormalTok{(}\DecValTok{2}\NormalTok{), }\AttributeTok{heading.title.font.size =} \StringTok{"15pt"}\NormalTok{,}
              \AttributeTok{heading.subtitle.font.size =} \StringTok{"11pt"}\NormalTok{, }\AttributeTok{latex.tbl.pos =} \StringTok{"H"}\NormalTok{)}
\end{Highlighting}
\end{Shaded}

\begin{table}
\caption*{
{\fontsize{15}{19}\selectfont  \textbf{Stationarity \& Differencing Diagnostic Tests}\fontsize{12}{15}\selectfont } \\ 
{\fontsize{11}{14}\selectfont  Fortified Wine Sales Training Period (1980–1993)\fontsize{12}{15}\selectfont }
} 
\fontsize{12.0pt}{14.0pt}\selectfont
\begin{tabular*}{\linewidth}{@{\extracolsep{\fill}}cccc}
\toprule
Diagnostic Test & Series Evaluated & Test Result & Conclusion \\ 
\midrule\addlinespace[2.5pt]
Seasonal Diff (nsdiffs) & Fortified (Raw) & D = 1 & 1 seasonal diff (D = 1) \\ 
First Diff (ndiffs) & Lag-12m diff & d = 0 & No regular diff (d = 0) \\ 
KPSS Test & Lag-12m diff & Stat: 0.0873 (p = 0.1) & Stationary (Fail to reject H0) \\ 
\bottomrule
\end{tabular*}
\end{table}

\textbf{What tests or reasoning did you use to determine the
differencing required?}

\begin{itemize}
\tightlist
\item
  First, I started with seasonal differencing because that is where the
  textbook recommends starting, and unitroot\_nsdiffs() recommended one
  seasonal difference. From there, I applied the seasonal difference at
  12 months to the raw data.
\item
  Second, I tested for a non-seasonal difference with unitroot\_ndiffs()
  and no non-seasonal differencing was required. unitroot\_kpss() could
  not reject the null hypothesis of stationary data (p-value
  \textgreater= 0.1).
\end{itemize}

\newpage

\begin{Shaded}
\begin{Highlighting}[]
\CommentTok{\# Plot differenced series and ACF/PACF plots }
\NormalTok{acf\_plot }\OtherTok{\textless{}{-}}\NormalTok{ trn\_d12 }\SpecialCharTok{|\textgreater{}} \FunctionTok{gg\_tsdisplay}\NormalTok{(d12, }\AttributeTok{plot\_type =} \StringTok{"partial"}\NormalTok{, }\AttributeTok{lag\_max =} \DecValTok{24}\NormalTok{)}

\CommentTok{\# Add helpful labels and adjust formatting}
\NormalTok{acf\_plot[[}\DecValTok{1}\NormalTok{]] }\OtherTok{\textless{}{-}}\NormalTok{ acf\_plot[[}\DecValTok{1}\NormalTok{]] }\SpecialCharTok{+}
  \FunctionTok{labs}\NormalTok{(}\AttributeTok{title =} \StringTok{"ACF \& PACF of Seasonally Differenced Fortified Sales [d=0, D=1, m=12]"}\NormalTok{) }\SpecialCharTok{+}
  \FunctionTok{theme}\NormalTok{(}\AttributeTok{plot.title =} \FunctionTok{element\_text}\NormalTok{(}\AttributeTok{size =} \DecValTok{11}\NormalTok{))}
\NormalTok{acf\_plot[[}\DecValTok{2}\NormalTok{]] }\OtherTok{\textless{}{-}}\NormalTok{ acf\_plot[[}\DecValTok{2}\NormalTok{]] }\SpecialCharTok{+}
  \FunctionTok{theme}\NormalTok{(}\AttributeTok{axis.text.x =} \FunctionTok{element\_text}\NormalTok{(}\AttributeTok{angle =} \DecValTok{45}\NormalTok{, }\AttributeTok{hjust =} \DecValTok{1}\NormalTok{, }\AttributeTok{size =} \DecValTok{8}\NormalTok{))}
\NormalTok{acf\_plot[[}\DecValTok{3}\NormalTok{]] }\OtherTok{\textless{}{-}}\NormalTok{ acf\_plot[[}\DecValTok{3}\NormalTok{]] }\SpecialCharTok{+}
  \FunctionTok{theme}\NormalTok{(}\AttributeTok{axis.text.x =} \FunctionTok{element\_text}\NormalTok{(}\AttributeTok{angle =} \DecValTok{45}\NormalTok{, }\AttributeTok{hjust =} \DecValTok{1}\NormalTok{, }\AttributeTok{size =} \DecValTok{8}\NormalTok{))}
\NormalTok{acf\_plot}
\end{Highlighting}
\end{Shaded}

\pandocbounded{\includegraphics[keepaspectratio]{ADS506-Assignment_3_1_files/figure-pdf/wine_fort_acf pt. 2-1.pdf}}

\textbf{What models do the ACF and PACF suggest would be best,
respectively?}

\begin{itemize}
\item
  ACF: The strong negative spike at lag 12 suggests a seasonal MA(1), no
  lags 1-11 are significant, meaning a non-seasonal MA value is not
  required. This suggests an ARIMA(0,0,0)(0,1,1){[}12{]}.
\item
  PACF: The strong negative spike at lag 12 suggests a seasonal AR(1),
  no lags 1-11 are significant, meaning a non-seasonal AR value is not
  required. This suggests an ARIMA(0,0,0)(1,1,0){[}12{]}.
\end{itemize}

\newpage

\subsubsection{b) Build the suggest models and an auto selected
model}\label{b-build-the-suggest-models-and-an-auto-selected-model}

\begin{Shaded}
\begin{Highlighting}[]
\CommentTok{\# Model fit}
\NormalTok{aus\_wine\_fit }\OtherTok{\textless{}{-}}\NormalTok{ trn }\SpecialCharTok{|\textgreater{}}
  \FunctionTok{model}\NormalTok{(}\AttributeTok{ma\_model =} \FunctionTok{ARIMA}\NormalTok{(Fortified }\SpecialCharTok{\textasciitilde{}} \FunctionTok{pdq}\NormalTok{(}\DecValTok{0}\NormalTok{, }\DecValTok{0}\NormalTok{, }\DecValTok{0}\NormalTok{) }\SpecialCharTok{+} \FunctionTok{PDQ}\NormalTok{(}\DecValTok{0}\NormalTok{,}\DecValTok{1}\NormalTok{,}\DecValTok{1}\NormalTok{)),}
        \AttributeTok{ar\_model =} \FunctionTok{ARIMA}\NormalTok{(Fortified }\SpecialCharTok{\textasciitilde{}} \FunctionTok{pdq}\NormalTok{(}\DecValTok{0}\NormalTok{, }\DecValTok{0}\NormalTok{, }\DecValTok{0}\NormalTok{) }\SpecialCharTok{+} \FunctionTok{PDQ}\NormalTok{(}\DecValTok{1}\NormalTok{,}\DecValTok{1}\NormalTok{,}\DecValTok{0}\NormalTok{)),}
        \AttributeTok{auto =} \FunctionTok{ARIMA}\NormalTok{(Fortified, }\AttributeTok{stepwise =} \ConstantTok{FALSE}\NormalTok{, }\AttributeTok{approx =} \ConstantTok{FALSE}\NormalTok{))}

\CommentTok{\# Report auto{-}arima model}
\NormalTok{aus\_wine\_fit }\SpecialCharTok{|\textgreater{}}
  \FunctionTok{pivot\_longer}\NormalTok{(}\FunctionTok{everything}\NormalTok{(), }\AttributeTok{names\_to =} \StringTok{"Model"}\NormalTok{, }\AttributeTok{values\_to =} \StringTok{"Specification"}\NormalTok{) }\SpecialCharTok{|\textgreater{}}
  \FunctionTok{as\_tibble}\NormalTok{() }\SpecialCharTok{|\textgreater{}}
  \FunctionTok{mutate}\NormalTok{(}\AttributeTok{Specification =} \FunctionTok{format}\NormalTok{(Specification)) }\SpecialCharTok{|\textgreater{}}
  \FunctionTok{filter}\NormalTok{(Model }\SpecialCharTok{==} \StringTok{"auto"}\NormalTok{) }\SpecialCharTok{|\textgreater{}}
  \FunctionTok{gt}\NormalTok{() }\SpecialCharTok{|\textgreater{}}
  \FunctionTok{tab\_header}\NormalTok{(}
    \AttributeTok{title =} \FunctionTok{md}\NormalTok{(}\StringTok{"**Fitted Auto ARIMA Model**"}\NormalTok{),}
    \AttributeTok{subtitle =} \StringTok{"Fortified Wine Sales (Training Period: 1980–1993)"}\NormalTok{) }\SpecialCharTok{|\textgreater{}}
  \FunctionTok{cols\_align}\NormalTok{(}\AttributeTok{align =} \StringTok{"center"}\NormalTok{, }\AttributeTok{columns =} \FunctionTok{everything}\NormalTok{()) }\SpecialCharTok{|\textgreater{}}
  \FunctionTok{tab\_options}\NormalTok{(}\AttributeTok{table.width =} \FunctionTok{pct}\NormalTok{(}\DecValTok{60}\NormalTok{), }\AttributeTok{latex.tbl.pos =} \StringTok{"H"}\NormalTok{)}
\end{Highlighting}
\end{Shaded}

\begin{table}
\caption*{
{\fontsize{20}{25}\selectfont  \textbf{Fitted Auto ARIMA Model}\fontsize{12}{15}\selectfont } \\ 
{\fontsize{14}{17}\selectfont  Fortified Wine Sales (Training Period: 1980–1993)\fontsize{12}{15}\selectfont }
} 
\fontsize{12.0pt}{14.0pt}\selectfont
\begin{tabular*}{0.6\linewidth}{@{\extracolsep{\fill}}cc}
\toprule
Model & Specification \\ 
\midrule\addlinespace[2.5pt]
auto & <ARIMA(0,0,1)(1,1,1)[12] w/ drift> \\ 
\bottomrule
\end{tabular*}
\end{table}

\textbf{What model did the ARIMA auto-select choose?}

The auto ARIMA model selected was ARIMA(0,0,1)(1,1,1){[}12{]} with a
drift term.

\newpage

\subsubsection{c) Compare the models with both RMSE and
AICc.}\label{c-compare-the-models-with-both-rmse-and-aicc.}

\begin{itemize}
\tightlist
\item
  Print the model metrics. Sort the report from best to worst based on
  \texttt{AICc}
\end{itemize}

\begin{Shaded}
\begin{Highlighting}[]
\CommentTok{\# Construct and view fit metrics table sorted by AICc}
\NormalTok{fit\_metrics }\OtherTok{\textless{}{-}} \FunctionTok{left\_join}\NormalTok{(}\FunctionTok{glance}\NormalTok{(aus\_wine\_fit) }\SpecialCharTok{|\textgreater{}}
  \FunctionTok{select}\NormalTok{(.model, AICc, AIC, BIC), }\FunctionTok{accuracy}\NormalTok{(aus\_wine\_fit) }\SpecialCharTok{|\textgreater{}}
  \FunctionTok{select}\NormalTok{(.model, RMSE, MAE, MAPE),}\AttributeTok{by =} \StringTok{".model"}\NormalTok{) }\SpecialCharTok{|\textgreater{}}
  \FunctionTok{rename}\NormalTok{(}\AttributeTok{Model =}\NormalTok{ .model) }\SpecialCharTok{|\textgreater{}} \FunctionTok{as\_tibble}\NormalTok{()}

\NormalTok{fit\_metrics }\SpecialCharTok{|\textgreater{}} \FunctionTok{arrange}\NormalTok{(AICc) }\SpecialCharTok{|\textgreater{}} \FunctionTok{gt}\NormalTok{() }\SpecialCharTok{|\textgreater{}}
  \FunctionTok{tab\_header}\NormalTok{(}\AttributeTok{title =} \FunctionTok{md}\NormalTok{(}\StringTok{"**Model Comparison: Information Criteria**"}\NormalTok{),}
             \AttributeTok{subtitle =} \StringTok{"Sorted by AICc (Training: 1980–1993)"}\NormalTok{) }\SpecialCharTok{|\textgreater{}}
  \FunctionTok{fmt\_number}\NormalTok{(}\AttributeTok{columns =} \FunctionTok{c}\NormalTok{(AICc, AIC, BIC, RMSE, MAE), }\AttributeTok{decimals =} \DecValTok{2}\NormalTok{) }\SpecialCharTok{|\textgreater{}}
  \FunctionTok{fmt\_percent}\NormalTok{(}\AttributeTok{columns =}\NormalTok{ MAPE, }\AttributeTok{scale\_values =} \ConstantTok{FALSE}\NormalTok{, }\AttributeTok{decimals =} \DecValTok{2}\NormalTok{) }\SpecialCharTok{|\textgreater{}}
  \FunctionTok{cols\_align}\NormalTok{(}\AttributeTok{align =} \StringTok{"center"}\NormalTok{, }\AttributeTok{columns =} \FunctionTok{everything}\NormalTok{()) }\SpecialCharTok{|\textgreater{}}
  \FunctionTok{tab\_options}\NormalTok{(}\AttributeTok{table.width =} \FunctionTok{pct}\NormalTok{(}\DecValTok{75}\NormalTok{), }\AttributeTok{latex.tbl.pos =} \StringTok{"H"}\NormalTok{)}
\end{Highlighting}
\end{Shaded}

\begin{table}
\caption*{
{\fontsize{20}{25}\selectfont  \textbf{Model Comparison: Information Criteria}\fontsize{12}{15}\selectfont } \\ 
{\fontsize{14}{17}\selectfont  Sorted by AICc (Training: 1980–1993)\fontsize{12}{15}\selectfont }
} 
\fontsize{12.0pt}{14.0pt}\selectfont
\begin{tabular*}{0.75\linewidth}{@{\extracolsep{\fill}}ccccccc}
\toprule
Model & AICc & AIC & BIC & RMSE & MAE & MAPE \\ 
\midrule\addlinespace[2.5pt]
auto & 2,230.99 & 2,230.59 & 2,245.84 & 280.59 & 207.48 & 6.80\% \\ 
ma\_model & 2,231.52 & 2,231.36 & 2,240.51 & 287.33 & 211.71 & 6.90\% \\ 
ar\_model & 2,241.23 & 2,241.07 & 2,250.22 & 299.26 & 223.69 & 7.38\% \\ 
\bottomrule
\end{tabular*}
\end{table}

\begin{itemize}
\tightlist
\item
  Print the model metrics. Sort the report from best to worst based on
  \texttt{RMSE}
\end{itemize}

\begin{Shaded}
\begin{Highlighting}[]
\CommentTok{\# View fit metrics table sorted by RMSE}
\NormalTok{fit\_metrics }\SpecialCharTok{|\textgreater{}} \FunctionTok{arrange}\NormalTok{(RMSE) }\SpecialCharTok{|\textgreater{}} \FunctionTok{gt}\NormalTok{() }\SpecialCharTok{|\textgreater{}}
  \FunctionTok{tab\_header}\NormalTok{(}\AttributeTok{title =} \FunctionTok{md}\NormalTok{(}\StringTok{"**Model Comparison: RMSE**"}\NormalTok{),}
             \AttributeTok{subtitle =} \StringTok{"Sorted by RMSE (Training: 1980–1993)"}\NormalTok{) }\SpecialCharTok{|\textgreater{}}
  \FunctionTok{fmt\_number}\NormalTok{(}\AttributeTok{columns =} \FunctionTok{c}\NormalTok{(AICc, AIC, BIC, RMSE, MAE), }\AttributeTok{decimals =} \DecValTok{2}\NormalTok{) }\SpecialCharTok{|\textgreater{}}
  \FunctionTok{fmt\_percent}\NormalTok{(}\AttributeTok{columns =}\NormalTok{ MAPE, }\AttributeTok{scale\_values =} \ConstantTok{FALSE}\NormalTok{, }\AttributeTok{decimals =} \DecValTok{2}\NormalTok{) }\SpecialCharTok{|\textgreater{}}
  \FunctionTok{cols\_align}\NormalTok{(}\AttributeTok{align =} \StringTok{"center"}\NormalTok{, }\AttributeTok{columns =} \FunctionTok{everything}\NormalTok{()) }\SpecialCharTok{|\textgreater{}}
  \FunctionTok{tab\_options}\NormalTok{(}\AttributeTok{table.width =} \FunctionTok{pct}\NormalTok{(}\DecValTok{75}\NormalTok{), }\AttributeTok{latex.tbl.pos =} \StringTok{"H"}\NormalTok{)}
\end{Highlighting}
\end{Shaded}

\begin{table}
\caption*{
{\fontsize{20}{25}\selectfont  \textbf{Model Comparison: RMSE}\fontsize{12}{15}\selectfont } \\ 
{\fontsize{14}{17}\selectfont  Sorted by RMSE (Training: 1980–1993)\fontsize{12}{15}\selectfont }
} 
\fontsize{12.0pt}{14.0pt}\selectfont
\begin{tabular*}{0.75\linewidth}{@{\extracolsep{\fill}}ccccccc}
\toprule
Model & AICc & AIC & BIC & RMSE & MAE & MAPE \\ 
\midrule\addlinespace[2.5pt]
auto & 2,230.99 & 2,230.59 & 2,245.84 & 280.59 & 207.48 & 6.80\% \\ 
ma\_model & 2,231.52 & 2,231.36 & 2,240.51 & 287.33 & 211.71 & 6.90\% \\ 
ar\_model & 2,241.23 & 2,241.07 & 2,250.22 & 299.26 & 223.69 & 7.38\% \\ 
\bottomrule
\end{tabular*}
\end{table}

\newpage

\subsubsection{d) Forecast the data for each month of 1994 for each
model and compare the resulting
metrics.}\label{d-forecast-the-data-for-each-month-of-1994-for-each-model-and-compare-the-resulting-metrics.}

\begin{itemize}
\tightlist
\item
  Report RMSE, MAE, and MAPE for all models, ordered by best RMSE.
\item
  Present your results in a well formatted table (suitable for
  presentation).
\end{itemize}

\begin{Shaded}
\begin{Highlighting}[]
\CommentTok{\# Forecast the validation period }
\NormalTok{forecast }\OtherTok{\textless{}{-}}\NormalTok{ aus\_wine\_fit }\SpecialCharTok{|\textgreater{}} \FunctionTok{forecast}\NormalTok{(}\AttributeTok{new\_data =}\NormalTok{ test)}

\CommentTok{\# Review accuracy to select the best model}
\FunctionTok{accuracy}\NormalTok{(forecast, test) }\SpecialCharTok{|\textgreater{}}
  \FunctionTok{select}\NormalTok{(}\AttributeTok{Model =}\NormalTok{ .model, RMSE, MAE, MAPE) }\SpecialCharTok{|\textgreater{}}
  \FunctionTok{arrange}\NormalTok{(RMSE) }\SpecialCharTok{|\textgreater{}}
  \FunctionTok{gt}\NormalTok{() }\SpecialCharTok{|\textgreater{}}
  \FunctionTok{tab\_header}\NormalTok{(}\AttributeTok{title =} \FunctionTok{md}\NormalTok{(}\StringTok{"**Model Comparison: Forecast Accuracy**"}\NormalTok{),}
    \AttributeTok{subtitle =} \StringTok{"Sorted from best to worst by RMSE (Testing: 1994)"}\NormalTok{) }\SpecialCharTok{|\textgreater{}}
  \FunctionTok{fmt\_number}\NormalTok{(}\AttributeTok{columns =} \FunctionTok{c}\NormalTok{(RMSE, MAE), }\AttributeTok{decimals =} \DecValTok{2}\NormalTok{) }\SpecialCharTok{|\textgreater{}}
  \FunctionTok{fmt\_percent}\NormalTok{(}\AttributeTok{columns =}\NormalTok{ MAPE, }\AttributeTok{scale\_values =} \ConstantTok{FALSE}\NormalTok{, }\AttributeTok{decimals =} \DecValTok{2}\NormalTok{) }\SpecialCharTok{|\textgreater{}}
  \FunctionTok{cols\_align}\NormalTok{(}\AttributeTok{align =} \StringTok{"center"}\NormalTok{, }\AttributeTok{columns =} \FunctionTok{everything}\NormalTok{()) }\SpecialCharTok{|\textgreater{}}
  \FunctionTok{tab\_options}\NormalTok{(}\AttributeTok{table.width =} \FunctionTok{pct}\NormalTok{(}\DecValTok{70}\NormalTok{), }\AttributeTok{latex.tbl.pos =} \StringTok{"H"}\NormalTok{)}
\end{Highlighting}
\end{Shaded}

\begin{table}
\caption*{
{\fontsize{20}{25}\selectfont  \textbf{Model Comparison: Forecast Accuracy}\fontsize{12}{15}\selectfont } \\ 
{\fontsize{14}{17}\selectfont  Sorted from best to worst by RMSE (Testing: 1994)\fontsize{12}{15}\selectfont }
} 
\fontsize{12.0pt}{14.0pt}\selectfont
\begin{tabular*}{0.7\linewidth}{@{\extracolsep{\fill}}cccc}
\toprule
Model & RMSE & MAE & MAPE \\ 
\midrule\addlinespace[2.5pt]
ar\_model & 285.89 & 218.31 & 9.31\% \\ 
ma\_model & 308.86 & 246.90 & 10.61\% \\ 
auto & 338.03 & 274.63 & 12.07\% \\ 
\bottomrule
\end{tabular*}
\end{table}

\textbf{Which model produces the best mean forecast? How did that
compare with the training goodness-of-fit metrics?}

The AR Model (ARIMA(0, 0, 0)(1,1,0)) shows the lowest out-of-sample RMSE
(285.89), MAE (218.31), and MAPE (9.31\%) of any models, followed by the
MA Model and the Auto model in last place. This is the exact opposite of
the training goodness of fit metrics where auto was in first place, MA
model in second, and AR in last place. The auto model clearly overfit
the training data, which is why the manual models outperformed the auto
model, highlighting the importance of configuring well thought out ARIMA
models using an ACF/PACF plot analysis.

\newpage

\subsubsection{e) Use the best performing model to create a forecast for
each month in
1995.}\label{e-use-the-best-performing-model-to-create-a-forecast-for-each-month-in-1995.}

\begin{itemize}
\tightlist
\item
  Fit a model against the full dataset.
\item
  Forecast the next year of data.
\item
  Create a plot showing forecast values and intervals (80\%, 95\%) for
  1995 alongside actual values from 1992 onwards.
\end{itemize}

\begin{Shaded}
\begin{Highlighting}[]
\CommentTok{\# Use the full data set to fit the best model}
\NormalTok{aus\_wine\_full\_fit }\OtherTok{\textless{}{-}}\NormalTok{ aus\_wine\_ts }\SpecialCharTok{|\textgreater{}}
  \FunctionTok{model}\NormalTok{(}\AttributeTok{ar\_model =} \FunctionTok{ARIMA}\NormalTok{(Fortified }\SpecialCharTok{\textasciitilde{}} \FunctionTok{pdq}\NormalTok{(}\DecValTok{0}\NormalTok{, }\DecValTok{0}\NormalTok{, }\DecValTok{0}\NormalTok{) }\SpecialCharTok{+} \FunctionTok{PDQ}\NormalTok{(}\DecValTok{1}\NormalTok{,}\DecValTok{1}\NormalTok{,}\DecValTok{0}\NormalTok{)))}

\CommentTok{\# Forecast the next 12 months}
\NormalTok{forecast }\OtherTok{\textless{}{-}}\NormalTok{ aus\_wine\_full\_fit }\SpecialCharTok{|\textgreater{}}
  \FunctionTok{forecast}\NormalTok{(}\AttributeTok{h=}\DecValTok{12}\NormalTok{)}

\CommentTok{\# Plot the forecast}
\NormalTok{forecast }\SpecialCharTok{|\textgreater{}}
  \FunctionTok{autoplot}\NormalTok{(aus\_wine\_ts }\SpecialCharTok{|\textgreater{}} \FunctionTok{filter}\NormalTok{(Month }\SpecialCharTok{\textgreater{}=} \FunctionTok{yearmonth}\NormalTok{(}\StringTok{"1992 Jan"}\NormalTok{))) }\SpecialCharTok{+}
  \FunctionTok{labs}\NormalTok{(}\AttributeTok{title =} \StringTok{"1995 Fortified Wine Sales Forecast"}\NormalTok{,}
       \AttributeTok{subtitle =} \StringTok{"ARIMA(0,0,0)(1,1,0)[12] with 80\% \& 95\% Prediction Intervals"}\NormalTok{,}
       \AttributeTok{x =} \StringTok{"Month"}\NormalTok{, }\AttributeTok{y =} \StringTok{"Sales (thousands of liters)"}\NormalTok{)}
\end{Highlighting}
\end{Shaded}

\pandocbounded{\includegraphics[keepaspectratio]{ADS506-Assignment_3_1_files/figure-pdf/arima-fc-plt-1.pdf}}

\newpage

\subsubsection{f) Output the mean forecast values for 1995 in a well
formatted
table.}\label{f-output-the-mean-forecast-values-for-1995-in-a-well-formatted-table.}

\begin{Shaded}
\begin{Highlighting}[]
\CommentTok{\# Output forecasted values by month}
\NormalTok{forecast }\SpecialCharTok{|\textgreater{}}
  \FunctionTok{as\_tibble}\NormalTok{() }\SpecialCharTok{|\textgreater{}}
  \FunctionTok{select}\NormalTok{(Month, }\AttributeTok{Mean =}\NormalTok{ .mean) }\SpecialCharTok{|\textgreater{}}
  \FunctionTok{gt}\NormalTok{() }\SpecialCharTok{|\textgreater{}}
  \FunctionTok{tab\_header}\NormalTok{(}\AttributeTok{title =} \FunctionTok{md}\NormalTok{(}\StringTok{"**1995 Mean Forecasted Fortified Wine Sales**"}\NormalTok{),}
    \AttributeTok{subtitle =} \StringTok{"Monthly Projections from Top{-}Performing ARIMA Model"}\NormalTok{) }\SpecialCharTok{|\textgreater{}}
  \FunctionTok{fmt\_number}\NormalTok{(}\AttributeTok{columns =}\NormalTok{ Mean, }\AttributeTok{decimals =} \DecValTok{2}\NormalTok{) }\SpecialCharTok{|\textgreater{}}
  \FunctionTok{cols\_align}\NormalTok{(}\AttributeTok{align =} \StringTok{"center"}\NormalTok{, }\AttributeTok{columns =} \FunctionTok{everything}\NormalTok{()) }\SpecialCharTok{|\textgreater{}}
  \FunctionTok{tab\_options}\NormalTok{(}\AttributeTok{table.width =} \FunctionTok{pct}\NormalTok{(}\DecValTok{50}\NormalTok{), }\AttributeTok{latex.tbl.pos =} \StringTok{"H"}\NormalTok{)}
\end{Highlighting}
\end{Shaded}

\begin{table}
\caption*{
{\fontsize{20}{25}\selectfont  \textbf{1995 Mean Forecasted Fortified Wine Sales}\fontsize{12}{15}\selectfont } \\ 
{\fontsize{14}{17}\selectfont  Monthly Projections from Top-Performing ARIMA Model\fontsize{12}{15}\selectfont }
} 
\fontsize{12.0pt}{14.0pt}\selectfont
\begin{tabular*}{0.5\linewidth}{@{\extracolsep{\fill}}cc}
\toprule
Month & Mean \\ 
\midrule\addlinespace[2.5pt]
1995 Jan & 1,013.97 \\ 
1995 Feb & 1,342.86 \\ 
1995 Mar & 1,802.68 \\ 
1995 Apr & 2,237.29 \\ 
1995 May & 2,181.95 \\ 
1995 Jun & 2,457.65 \\ 
1995 Jul & 2,636.80 \\ 
1995 Aug & 2,266.65 \\ 
1995 Sep & 2,138.98 \\ 
1995 Oct & 1,755.19 \\ 
1995 Nov & 2,526.44 \\ 
1995 Dec & 2,421.81 \\ 
\bottomrule
\end{tabular*}
\end{table}

\begin{center}
\textit{Note: Values reported in thousands of liters.}
\end{center}

\newpage

\subsection{2) Forecasting Australian Wine Sales: ARIMA comparison
against other
models}\label{forecasting-australian-wine-sales-arima-comparison-against-other-models}

\begin{itemize}
\tightlist
\item
  Use your existing partitioned data for the fortified wine series
  (training through December 1993).
\item
  Fit an ARIMA model - use the best model from the previous problem
\item
  Also fit time series linear model with trend and seasonal components
  and an auto-selected exponential smoothing model. (Hint: see the
  assignment from the previous module).
\end{itemize}

\subsubsection{a) Create the ``actual vs.~forecast'' plot for validation
period.}\label{a-create-the-actual-vs.-forecast-plot-for-validation-period.}

\begin{itemize}
\tightlist
\item
  (Hint) It may be easier to create a new \texttt{mable} object by
  refitting the data
\item
  plot all 3 models together with the actual data
\item
  remove the confidence intervals for clarity
\end{itemize}

\begin{Shaded}
\begin{Highlighting}[]
\CommentTok{\# Build models and forecast validation period}
\NormalTok{aus\_wine\_fit }\OtherTok{\textless{}{-}}\NormalTok{ trn }\SpecialCharTok{|\textgreater{}}
  \FunctionTok{model}\NormalTok{(}\AttributeTok{ar\_model =} \FunctionTok{ARIMA}\NormalTok{(Fortified }\SpecialCharTok{\textasciitilde{}} \FunctionTok{pdq}\NormalTok{(}\DecValTok{0}\NormalTok{, }\DecValTok{0}\NormalTok{, }\DecValTok{0}\NormalTok{) }\SpecialCharTok{+} \FunctionTok{PDQ}\NormalTok{(}\DecValTok{1}\NormalTok{,}\DecValTok{1}\NormalTok{,}\DecValTok{0}\NormalTok{)),}
        \AttributeTok{tslm =} \FunctionTok{TSLM}\NormalTok{(Fortified }\SpecialCharTok{\textasciitilde{}} \FunctionTok{trend}\NormalTok{() }\SpecialCharTok{+} \FunctionTok{season}\NormalTok{()),}
        \AttributeTok{ets =} \FunctionTok{ETS}\NormalTok{(Fortified))}
\NormalTok{forecast }\OtherTok{\textless{}{-}}\NormalTok{ aus\_wine\_fit }\SpecialCharTok{|\textgreater{}} \FunctionTok{forecast}\NormalTok{(}\AttributeTok{h =} \DecValTok{12}\NormalTok{)}

\CommentTok{\# Plot validation forecasts against actuals}
\NormalTok{forecast }\SpecialCharTok{|\textgreater{}}
  \FunctionTok{autoplot}\NormalTok{(aus\_wine\_ts }\SpecialCharTok{|\textgreater{}} \FunctionTok{filter}\NormalTok{(Month }\SpecialCharTok{\textgreater{}=} \FunctionTok{yearmonth}\NormalTok{(}\StringTok{"1992 Jan"}\NormalTok{)), }\AttributeTok{level =} \ConstantTok{NULL}\NormalTok{) }\SpecialCharTok{+}
  \FunctionTok{labs}\NormalTok{(}\AttributeTok{title =} \StringTok{"1994 Validation Forecast Comparison"}\NormalTok{,}
       \AttributeTok{subtitle =} \StringTok{"Comparing AR model, ETS, and TSLM Models Against Actual Sales"}\NormalTok{,}
       \AttributeTok{x =} \StringTok{"Month"}\NormalTok{, }\AttributeTok{y =} \StringTok{"Sales (thousands of liters)"}\NormalTok{)}
\end{Highlighting}
\end{Shaded}

\pandocbounded{\includegraphics[keepaspectratio]{ADS506-Assignment_3_1_files/figure-pdf/fortified-compare-1.pdf}}

\newpage

\begin{itemize}
\tightlist
\item
  Compare the RSME, MAE and MAPE metrics in the validation period for
  each model in a well formatted table (suitable for presentation). Sort
  from best to worst RMSE.
\end{itemize}

\begin{Shaded}
\begin{Highlighting}[]
\CommentTok{\# View forecast accuracy table}
\FunctionTok{accuracy}\NormalTok{(forecast, test) }\SpecialCharTok{|\textgreater{}}
  \FunctionTok{select}\NormalTok{(}\AttributeTok{Model =}\NormalTok{ .model, RMSE, MAE, MAPE) }\SpecialCharTok{|\textgreater{}}
  \FunctionTok{arrange}\NormalTok{(RMSE) }\SpecialCharTok{|\textgreater{}}
  \FunctionTok{gt}\NormalTok{() }\SpecialCharTok{|\textgreater{}}
  \FunctionTok{tab\_header}\NormalTok{(}\AttributeTok{title =} \FunctionTok{md}\NormalTok{(}\StringTok{"**Model Comparison: Forecast Accuracy**"}\NormalTok{),}
    \AttributeTok{subtitle =} \StringTok{"Sorted from best to worst by RMSE (Testing: 1994)"}\NormalTok{) }\SpecialCharTok{|\textgreater{}}
  \FunctionTok{fmt\_number}\NormalTok{(}\AttributeTok{columns =} \FunctionTok{c}\NormalTok{(RMSE, MAE), }\AttributeTok{decimals =} \DecValTok{2}\NormalTok{) }\SpecialCharTok{|\textgreater{}}
  \FunctionTok{fmt\_percent}\NormalTok{(}\AttributeTok{columns =}\NormalTok{ MAPE, }\AttributeTok{scale\_values =} \ConstantTok{FALSE}\NormalTok{, }\AttributeTok{decimals =} \DecValTok{2}\NormalTok{) }\SpecialCharTok{|\textgreater{}}
  \FunctionTok{cols\_align}\NormalTok{(}\AttributeTok{align =} \StringTok{"center"}\NormalTok{, }\AttributeTok{columns =} \FunctionTok{everything}\NormalTok{()) }\SpecialCharTok{|\textgreater{}}
  \FunctionTok{tab\_options}\NormalTok{(}\AttributeTok{table.width =} \FunctionTok{pct}\NormalTok{(}\DecValTok{70}\NormalTok{), }\AttributeTok{latex.tbl.pos =} \StringTok{"H"}\NormalTok{)}
\end{Highlighting}
\end{Shaded}

\begin{table}
\caption*{
{\fontsize{20}{25}\selectfont  \textbf{Model Comparison: Forecast Accuracy}\fontsize{12}{15}\selectfont } \\ 
{\fontsize{14}{17}\selectfont  Sorted from best to worst by RMSE (Testing: 1994)\fontsize{12}{15}\selectfont }
} 
\fontsize{12.0pt}{14.0pt}\selectfont
\begin{tabular*}{0.7\linewidth}{@{\extracolsep{\fill}}cccc}
\toprule
Model & RMSE & MAE & MAPE \\ 
\midrule\addlinespace[2.5pt]
ar\_model & 285.89 & 218.31 & 9.31\% \\ 
ets & 318.90 & 242.65 & 10.24\% \\ 
tslm & 426.53 & 362.69 & 16.25\% \\ 
\bottomrule
\end{tabular*}
\end{table}

\textbf{\emph{Which model produced the best forecast based on RMSE?}}

The AR ARIMA model outperformed the ETS and TSLM models with a 285.89
RMSE, 218.31 MAE, and a 9.31\% MAPE.

\newpage

\subsubsection{b) Produce a forecast for all wine types for sales in all
months of 1995 using an auto ARIMA
model.}\label{b-produce-a-forecast-for-all-wine-types-for-sales-in-all-months-of-1995-using-an-auto-arima-model.}

\begin{itemize}
\tightlist
\item
  Output a presentable table with wine varietals as columns and months
  as rows
\end{itemize}

\begin{Shaded}
\begin{Highlighting}[]
\CommentTok{\# Create a tsibble for all the wine types}
\NormalTok{all\_aus\_wine\_ts }\OtherTok{\textless{}{-}}\NormalTok{ aus\_wine }\SpecialCharTok{|\textgreater{}}
  \FunctionTok{pivot\_longer}\NormalTok{(}\AttributeTok{cols =} \SpecialCharTok{{-}}\NormalTok{Month, }\AttributeTok{names\_to =} \StringTok{\textquotesingle{}Varietal\textquotesingle{}}\NormalTok{, }\AttributeTok{values\_to =} \StringTok{\textquotesingle{}Sales\textquotesingle{}}\NormalTok{) }\SpecialCharTok{|\textgreater{}}
  \FunctionTok{as\_tsibble}\NormalTok{(}\AttributeTok{index =}\NormalTok{ Month, }\AttributeTok{key =} \StringTok{\textquotesingle{}Varietal\textquotesingle{}}\NormalTok{)}

\CommentTok{\# Fit auto{-}arima}
\NormalTok{all\_aus\_wine\_fit }\OtherTok{\textless{}{-}}\NormalTok{ all\_aus\_wine\_ts }\SpecialCharTok{|\textgreater{}} \FunctionTok{model}\NormalTok{(}\AttributeTok{auto\_arima =} \FunctionTok{ARIMA}\NormalTok{(Sales))}

\CommentTok{\# Forecast}
\NormalTok{forecast }\OtherTok{\textless{}{-}}\NormalTok{ all\_aus\_wine\_fit }\SpecialCharTok{|\textgreater{}} \FunctionTok{forecast}\NormalTok{(}\AttributeTok{h =} \DecValTok{12}\NormalTok{)}

\CommentTok{\# Use pivot wider to shape the output}
\NormalTok{forecast }\SpecialCharTok{|\textgreater{}}
  \FunctionTok{as\_tibble}\NormalTok{() }\SpecialCharTok{|\textgreater{}}
  \FunctionTok{select}\NormalTok{(Varietal, Month, }\AttributeTok{Mean =}\NormalTok{ .mean) }\SpecialCharTok{|\textgreater{}}
  \FunctionTok{pivot\_wider}\NormalTok{(}\AttributeTok{names\_from =}\NormalTok{ Varietal, }\AttributeTok{values\_from =}\NormalTok{ Mean) }\SpecialCharTok{|\textgreater{}}
  \FunctionTok{rename}\NormalTok{(}\StringTok{\textasciigrave{}}\AttributeTok{Dry White}\StringTok{\textasciigrave{}} \OtherTok{=} \StringTok{\textasciigrave{}}\AttributeTok{Dry white}\StringTok{\textasciigrave{}}\NormalTok{, }\StringTok{\textasciigrave{}}\AttributeTok{Sweet White}\StringTok{\textasciigrave{}} \OtherTok{=} \StringTok{\textasciigrave{}}\AttributeTok{Sweet white}\StringTok{\textasciigrave{}}\NormalTok{, }\AttributeTok{Sparkling =}\NormalTok{ sparkling) }\SpecialCharTok{|\textgreater{}}
  \FunctionTok{gt}\NormalTok{() }\SpecialCharTok{|\textgreater{}}
  \FunctionTok{tab\_header}\NormalTok{(}\AttributeTok{title =} \FunctionTok{md}\NormalTok{(}\StringTok{"**1995 Auto{-}ARIMA Sales Forecast by Varietal**"}\NormalTok{),}
    \AttributeTok{subtitle =} \StringTok{"Monthly Mean Forecasts (Thousands of Liters)"}\NormalTok{) }\SpecialCharTok{|\textgreater{}}
  \FunctionTok{fmt\_number}\NormalTok{(}\AttributeTok{decimals =} \DecValTok{2}\NormalTok{) }\SpecialCharTok{|\textgreater{}}
  \FunctionTok{cols\_align}\NormalTok{(}\AttributeTok{align =} \StringTok{"center"}\NormalTok{, }\AttributeTok{columns =} \FunctionTok{everything}\NormalTok{()) }\SpecialCharTok{|\textgreater{}}
  \FunctionTok{tab\_options}\NormalTok{(}\AttributeTok{table.width =} \FunctionTok{pct}\NormalTok{(}\DecValTok{100}\NormalTok{), }\AttributeTok{latex.tbl.pos =} \StringTok{"H"}\NormalTok{)}
\end{Highlighting}
\end{Shaded}

\begin{table}
\caption*{
{\fontsize{20}{25}\selectfont  \textbf{1995 Auto-ARIMA Sales Forecast by Varietal}\fontsize{12}{15}\selectfont } \\ 
{\fontsize{14}{17}\selectfont  Monthly Mean Forecasts (Thousands of Liters)\fontsize{12}{15}\selectfont }
} 
\fontsize{12.0pt}{14.0pt}\selectfont
\begin{tabular*}{1\linewidth}{@{\extracolsep{\fill}}ccccccc}
\toprule
Month & Dry White & Fortified & Red & Rose & Sweet White & Sparkling \\ 
\midrule\addlinespace[2.5pt]
1995 Jan & 2,683.51 & 992.20 & 1,337.37 & 15.93 & 153.33 & 1,381.45 \\ 
1995 Feb & 3,685.31 & 1,317.51 & 1,795.36 & 27.68 & 218.29 & 1,788.24 \\ 
1995 Mar & 3,850.03 & 1,778.32 & 2,290.58 & 33.01 & 192.44 & 1,826.06 \\ 
1995 Apr & 3,714.70 & 2,189.91 & 2,412.17 & 30.29 & 221.36 & 1,829.86 \\ 
1995 May & 3,478.50 & 2,263.47 & 2,414.95 & 30.72 & 179.09 & 1,707.46 \\ 
1995 Jun & 3,460.96 & 2,454.42 & 2,640.51 & 36.38 & 176.70 & 1,608.55 \\ 
1995 Jul & 4,062.53 & 2,783.23 & 3,421.37 & 42.49 & 232.49 & 2,049.57 \\ 
1995 Aug & 4,074.40 & 2,351.13 & 2,959.91 & 39.05 & 228.14 & 1,909.54 \\ 
1995 Sep & 3,924.65 & 2,073.04 & 2,565.39 & 38.44 & 231.61 & 2,500.64 \\ 
1995 Oct & 4,183.30 & 1,760.84 & 2,347.85 & 39.75 & 243.44 & 3,321.33 \\ 
1995 Nov & 4,860.60 & 2,428.03 & 2,655.92 & 50.01 & 291.07 & 3,952.61 \\ 
1995 Dec & 5,634.42 & 2,410.22 & 2,809.16 & 81.20 & 390.58 & 6,157.06 \\ 
\bottomrule
\end{tabular*}
\end{table}

\newpage

\begin{itemize}
\tightlist
\item
  Plot the forecast for each varietal along with 3 years of historical
  data.
\end{itemize}

\begin{Shaded}
\begin{Highlighting}[]
\CommentTok{\# Show forecast plots}
\NormalTok{forecast }\SpecialCharTok{|\textgreater{}}
  \FunctionTok{autoplot}\NormalTok{(all\_aus\_wine\_ts }\SpecialCharTok{|\textgreater{}} \FunctionTok{filter}\NormalTok{(}\FunctionTok{year}\NormalTok{(Month) }\SpecialCharTok{\textgreater{}=} \DecValTok{1992}\NormalTok{)) }\SpecialCharTok{+}
  \FunctionTok{facet\_wrap}\NormalTok{(}\SpecialCharTok{\textasciitilde{}}\NormalTok{ Varietal, }\AttributeTok{scales =} \StringTok{"free\_y"}\NormalTok{, }\AttributeTok{ncol =} \DecValTok{3}\NormalTok{) }\SpecialCharTok{+}
  \FunctionTok{labs}\NormalTok{(}\AttributeTok{title =} \StringTok{"1995 Sales Forecast by Wine Varietal"}\NormalTok{,}
    \AttributeTok{subtitle =} \StringTok{"Auto{-}ARIMA 12{-}Month Projections with 3 Years Historical Data (1992–1994)"}\NormalTok{,}
    \AttributeTok{x =} \StringTok{"Month"}\NormalTok{, }\AttributeTok{y =} \StringTok{"Sales (thousands of liters)"}\NormalTok{) }\SpecialCharTok{+}
  \FunctionTok{theme}\NormalTok{(}\AttributeTok{legend.position =} \StringTok{"bottom"}\NormalTok{, }\AttributeTok{axis.text.x =} \FunctionTok{element\_text}\NormalTok{(}\AttributeTok{angle =} \DecValTok{45}\NormalTok{, }\AttributeTok{hjust =} \DecValTok{1}\NormalTok{))}
\end{Highlighting}
\end{Shaded}

\pandocbounded{\includegraphics[keepaspectratio]{ADS506-Assignment_3_1_files/figure-pdf/unnamed-chunk-2-1.pdf}}

\newpage

\subsection{3) Reflections}\label{reflections}

a. What was the most challenging part of this assignment?

\begin{itemize}
\tightlist
\item
  The most challenging part of this assignment was making sure the ACF
  and PACF analysis for what pdq and PDQ values to select was correct. I
  found the text and office hours lecture very helpful to confirm my
  thoughts here, but I did carefully have to consider that question.
\end{itemize}

b. What was your most valuable takeaway from this assignment?

\begin{itemize}
\tightlist
\item
  Manual ARIMA models often outperform auto ARIMA models out-of-sample.
  I am no stranger to the concept of overfitting, but it seems like auto
  ARIMA overfits often, which means that manually analyzing ACF and PACF
  plots is critical skill.
\end{itemize}

c. What improvements do you suggest for future courses?

\begin{itemize}
\tightlist
\item
  No improvements for future courses. I think it was great that TSLM and
  ETS were reinforced here, I think its important to keep practicing the
  syntax and understanding the outputs.
\end{itemize}

d. What LLMs did you use (if any) to help you complete this assignment?

\begin{itemize}
\tightlist
\item
  I used Gemini to assist me with speeding up great tables formatting
  and providing feedback on problems after I took a first and second
  pass over them. I consulted the text first to ensure I did not rely
  too heavily on Gemini.
\end{itemize}

\textbf{Reference}

Hyndman, R.J., \& Athanasopoulos, G. (2021) \emph{Forecasting:
principles and practice}, 3rd edition, OTexts: Melbourne, Australia.
OTexts.com/fpp3. Accessed on 30 October 2026.

\textbf{AI Disclosure:}

I used \textbf{Antigravity (powered by Gemini)} as a pair programmer for
technical guidance during this assignment. It assisted with debugging R
syntax. All code was executed, tested, and verified locally by me, and
all written responses reflect my own understanding and phrasing.




\end{document}
